\section{Background}
\label{sec:background}

\subsection{Satisfiability Modulo Theories and the SMT-LIB2 Standard}
\wyc{no...we never explain challenge in background section. Maybe you can explain the basic syntax of SMT-lib2, since it is used in the example section.}
\SMT{} is a fundamental decision problem for logical formulas that extends \SAT{} with a combination of background theories, such as arithmetic, bit-vectors, and arrays~\cite{kroening2016decision}. This ability to reason about expressive, domain-specific constraints has made \SMT{} a cornerstone of modern computer science, providing the core reasoning engine for diverse applications ranging from software verification and model checking~\cite{ball2004slam} to automated test generation via symbolic execution~\cite{cadar2008klee}. The widespread success and practical impact of SMT have spurred the development of numerous highly-optimized solvers~\cite{demoura2008z3,mathsat5,cvc5}.

AS a diverse ecosystem of powerful \SMT{} solvers emerged, each with distinct specializations and performance characteristics, a critical need for a common interchange format arose. The SMT-LIB2 format~\cite{barrett2010smtlib} was established as the community's solution, providing a unified and solver-agnostic standard. Its syntax is based on S-expressions, using a Lisp-like prefix notation where an operator precedes its arguments (\eg{}, \texttt{(+ x 5)} represents the expression $x+5$). A problem definition in SMT-LIB2 is typically structured using a few core commands.
The \texttt{declare-const} command, written as \texttt{(declare-const <symbol> <sort>)}, introduces a symbolic constant (\ie{}, variable) of a specified sort into the solver's logical context. For example, \texttt{(declare-const x Int)} declares a symbolic integer named \texttt{x} that can be referenced throughout the formula. More generally, the \texttt{declare-fun} command, written as \texttt{(declare-fun <symbol> (<sorts>) <sort>)}, declares a function symbol, where \texttt{declare-const} serves as a convenient shorthand for declaring a function with no arguments.
To manage complexity in formulas with deeply nested expressions, such as those presented in our motivating example (Section~\ref{sec:motivation-and-problem}), the \texttt{let} construct, written as \texttt{(let ([<bindings>]) <term>)}, is essential. This construct allows one or more complex sub-expressions to be assigned to local variables within the \texttt{[<bindings>]} block for use within the body \texttt{<term>}, which significantly improves a formula's readability and structure by avoiding repetitive sub-expressions.
The \texttt{assert} command, written as \texttt{(assert <formula>)}, introduces a logical constraint---namely, the Boolean formula \texttt{<formula>}---into the solver's current assertion stack. This stack accumulates all active constraints, which are collectively evaluated when a satisfiability check is performed. The solver's primary goal is to determine if there exists a valid assignment for all declared constants that makes all asserted constraints simultaneously true, for which it will return one of three results: \texttt{sat} (satisfiable), \texttt{unsat} (unsatisfiable), or \texttt{unknown} (typically due to timeout or resource limits).

\subsection{Reinforcement Learning for Sequential Decision Making}
\wyc{We discuss Reinforcement Learning for Sequential Decision Making. No need to discuss challenges in SMT query.

The outline could be: 
1. What is Reinforcement Learning, the basic framework of RL (one or two paragraphs)
2. How to apply RL to a downstream task (one paragraph).
}
\wyc{What is MDP? It is used in the later section but never explained. }
\lz{introduce RL and MDP}
\RL{} is a computational approach that addresses sequential decision-making problems by formalizing them as an interaction between two main components: a learning agent and its environment~\cite{sutton2018reinforcement}. The agent learns to achieve a goal through a process of trial-and-error, where at each time step it observes the environment's state, selects an action, and receives a corresponding reward and the next state as feedback. In this paradigm, the environment is formally modeled as a \MDP{}~\cite{puterman2014markov}, which serves as the complete mathematical specification of the learning problem. 
A standard \MDP{} is defined by a 4-tuple $(\mathcal{S}, \mathcal{A}, P, R)$, where $\mathcal{S}$ represents the set of all possible states, $\mathcal{A}$ represents the available actions, $P$ is the state transition dynamics, and $R$ is the reward function. The agent's fundamental task is then to learn a policy ($\pi: \mathcal{S} \rightarrow \mathcal{A}$)---a mapping from states to actions---that maximizes the cumulative reward it can collect.

For tasks with naturally bounded episodes, a finite-horizon \MDP{} extends this formulation with an additional component, yielding a 5-tuple $(\mathcal{S}, \mathcal{A}, P, R, H)$, where $H$ represents the finite horizon---the maximum number of time steps in each episode~\cite{sutton2018reinforcement}. This finite-horizon constraint is particularly relevant for episodic tasks where each problem instance has a natural termination condition, such as reaching a goal state or exhausting available resources. In finite-horizon settings, the agent's objective becomes maximizing the expected cumulative reward within the bounded time horizon $H$.

\lz{explain how to apply RL to a real-world task}
Before an agent can learn such a policy to deal with a real-world task, the task's characteristics must be translated into the formal components of an \MDP{}. This translation involves three key design choices: defining a state space ($\mathcal{S}$) that effectively captures all features of the task relevant for decision-making; defining an action space ($\mathcal{A}$) that encompasses all valid moves the agent can make; and designing a reward function ($R$) that provides feedback correctly aligned with the task's ultimate goal. For instance, in a classic grid-world navigation problem, an agent's state is typically its ($x, y$) coordinate, its actions are the set `\{up, down, left, right\}', and its reward function might provide $+1$ for reaching a target and $-1$ for hitting an obstacle. Once a problem is framed in this way, a suitable RL algorithm can be applied, allowing the agent to learn a policy that navigates the grid efficiently.

\lz{explain the challenges of applying RL to a real-world task, particularly the reward function, correspond with the reward in methods}
While the MDP framework is conceptually straightforward, designing its components for complex tasks presents significant challenges, particularly the reward function. Feedback from the environment is often naturally sparse---meaning a non-zero reward is received only after a long sequence of actions. This sparsity leads to the difficult temporal credit assignment problem, making learning inefficient~\cite{sutton2018reinforcement}. Algorithmic approaches such as Actor-Critic methods are designed to mitigate this issue~\cite{mnih2016asynchronous}. They introduce a critic that learns a long-term value function, providing the actor (the policy) with a denser, internally-generated estimate of future rewards to guide its learning. However, since the critic's estimates are themselves ultimately derived from the sparse environmental feedback, this internal guidance can be slow to converge and unreliable in highly complex settings, a well-known challenge in \RL{}~\cite{fujimoto2018addressing}. Consequently, effectively addressing severe reward sparsity often remains a fundamental problem in applied \RL{}, frequently necessitating the design of problem-specific, hybrid reward systems that incorporate multiple, diverse signal sources~\cite{akkaya2019solving}.

\subsection{Large Language Models for Semantic Reasoning}
\wyc{Usually introduce basic information about large language models.}
\LLM{}s are deep neural networks, typically based on the Transformer architecture~\cite{vaswani2017attention}, that are pre-trained on vast amounts of text and code~\cite{devlin2019bert}. This extensive pre-training equips them with a powerful internal model of language, encompassing syntax, semantics, and common-sense knowledge. The primary operational paradigm for modern LLMs is \textit{in-context learning}~\cite{brown2020language}, where they can be steered to perform novel tasks based solely on the context provided in a carefully constructed prompt, without any changes to the model's parameters.

Beyond fluency in natural language, this pre-training has led to surprising \textit{emergent abilities}~\cite{wei2022emergent}, with recent work demonstrating \LLM{}s' growing capacity for formal and heuristic reasoning in tasks such as invariant generation and logical problem solving~\cite{kamath2023finding, pan2023logiclm}. A key manifestation of this semantic reasoning is semantic synthesis---the ability to not just understand a context, but to actively generate a plausible continuation or structured output that adheres to its semantic constraints. It is this capability that allows an \LLM{} to function as a powerful engine for proposing candidate solutions that are consistent with a given problem's requirements.


